% Diese Datei erzeugt die typische Uni-Wuppertal Titelseite bestehend 
% aus dem gelben Streifen mit dem Bergischen Löwen,
% dem Schriftzug "Fakultät ..." sowie Namen und Titeln



\begin {titlepage}

\unitlength1mm
\begin{picture}(150,205)(0,40)


\put(0,150){\includegraphics[height=10cm]{boxleo}}

% Trick: 5cm hohe Box mit linken umteren Punkt 5cm unterhalb 
% von Oberkante des Streifens => Oberkante text = Oberkante Balken

% Unabhängigigkeit von UTF8 oder LATIN: Universit"at
\put(25,200){\parbox[b][5cm][t]{20cm}{\textsc{\large 
Bergische Universit\"at Wuppertal \\
Fakultät für \\
{\it Mathematik und  Naturwissenschaften} \\ 
Fachgruppe Physik \\ 
~\\
Fortgeschrittenen Praktikum}}}




\put (25,185.0){\parbox[b]{20cm}{ \Large \bf \titel}}



\put (25,152) {\parbox[b]{20cm}{\large \bf \autoren}}

\put (25,125) {\parbox[t]{11cm}{\large \bf Abstract (Kurzbeschreibung) \\
                                \normalsize\rm
                                \abstracttext}}

\put (25,60) {Abgabedatum: \datum}
\put (25,65) {Betreuer: \betreuer}



\put (22,30)
{
%  \fbox{
%    \begin{turn}{-90}
    \begin{tabular}{|l|c|l|c|}
    \hline
    Bewertung Protokoll   & max. \% & +/0/- &  erreicht  \% \\  \hline \hline
    Formales              & 6      &       & \\ \hline
    Einleitung \& Theorie &  6     &       & \\  \hline        
    Durchf\"uhrung        &        &       & \\
    Auswertung		      & 33     &       & \\ 
    phys. Diskussion      &        &       & \\ 
    Zusammenfassung       &        &       & \\ \hline \hline
    Qualit"at der Messung & 15     &       & \\ \hline
    \end{tabular}
    
%    \end{turn}
%  }
}

\end {picture}
\end {titlepage}








